\chapter*{Abstract}
\addcontentsline{toc}{chapter}{Abstract}

Multichannel and multistatic synthetic aperture radar (SAR) systems for high-resolution wide-swath (HRWS) imaging sample
the azimuth signal of each receiver at a pulse repetition frequency (PRF) below the Doppler bandwidth and recover the
unambiguous spectrum by combining the channels in an azimuth reconstruction. The reconstruction filters are derived for
a reference geometry. When the receivers are separated by cross-track baselines, a scatterer at a height different from
the reference height carries a channel-dependent phase, the aliased parts of the spectrum no longer cancel, and azimuth
ambiguities appear.

This report presents a \emph{joint} topography correction and azimuth reconstruction that removes these ambiguities using
\emph{only the system parameters and a digital elevation model (DEM)}, never the positions of the targets. The phase that
the terrain adds to each channel is computed from the DEM and placed inside the model of the channels, at every position
of the line; the focused line is then estimated by regularised least squares with conjugate gradients, applying the model
and its adjoint with FFTs only. The method is developed and tested with a numerical simulator of a spaceborne formation
($\lambda=\SI{0.25}{m}$, $r_0=\SI{766}{km}$, $\Nrx=2$ to $8$ receivers, cross-track baselines up to \SI{500}{m}).

Over \SweepOneRuns{} one-dimensional runs (\SweepOneScenes{} scenes, $\Nrx=2$--$8$) the median highest ambiguity drops
from \SI{\SweepOneMedNone}{dB} without correction to \SI{\SweepOneMedJoint}{dB} with the joint reconstruction
(reference \SI{\SweepOneMedRef}{dB}). The two-dimensional extension, which applies the DEM at every pixel of a 3-D scene,
reaches a median of \SI{\SweepTwoMedJoint}{dB} over \SweepTwoRuns{} runs (no correction \SI{\SweepTwoMedNone}{dB},
reference \SI{\SweepTwoMedRef}{dB}). On real terrain, 200 azimuth lines cut from the SRTM DEM of the Sierra Nevada, the
median goes from \SI{-8.3}{dB} to \SI{-41.5}{dB}, at the level of the reference (\SI{-41.6}{dB}). The method works best
with three or more channels and cross-track baselines that are not equally spaced, which give the channels the diversity
needed to separate the folded parts of the spectrum.
